% KT74011C
% Kindesmisshandlung
% 14.0
% 2013-11-30
\label{m:kindesmisshandlung}
Die Aufgaben des Rettungsdienstes bei Verdacht auf Kindesmisshandlung besteht in:

\begin{enumerate}
  \item Daran denken!
  \item Umstände und Umfeld abklären
  \item Unstimmigkeiten und Auffälligkeiten auf eine neutrale Weise
  \EE{dokumentieren}
  \item Obiges verlässlich bei der Übergabe an die nachbetreuende Einrichtung
  weiterleiten. Evtl. Name des Gesprächpartners bei der  Übergabe dokumentieren.
  \item \EE{Patienten nicht im Stich lassen}: Gelindestes Mittel wählen,
  evtl. Vorwand für die Hospitalisierung suchen, aber wenn ein begründeter
  Verdacht auf eine Misshandlung besteht und die
  Erziehungsberechtigten unkooperativ sind notfalls auch die Exekutive beiziehen.
\end{enumerate}

\paragraph{Vorsicht!}

Jedenfalls zu \underline{unterlassen} sind:

\begin{itemize}
  \item Voreilige Verdachtsäußerungen. \Z{Nicht alles was stinkt ist ein
  Fisch!}
  \item Andeutungen gegenüber den Erziehungsberechtigten oder Dritten.
\end{itemize}
